\section{Unique Paths II (With Obstacles)}
\label{sec:dp-unique-paths-2}

\problemstatement{Same as Section~\ref{sec:dp-unique-paths}, except some
grid cells are marked as \textbf{obstacles} (value $1$; open cells are
$0$), and the robot cannot step onto an obstacle cell at all. Count the
number of unique paths from top-left to bottom-right.}

\begin{patternbox}
\emph{``Same grid-path-counting problem, with some cells forbidden''}
needs exactly one change to Section~\ref{sec:dp-unique-paths}'s
recurrence: an obstacle cell contributes \textbf{zero} paths, regardless
of what its predecessors computed, because no path is allowed to pass
through it at all.
\end{patternbox}

\subsection*{Deriving the recurrence}

\[
  \textit{ways}(i,j) =
  \begin{cases}
    0 & \text{if } (i,j) \text{ is an obstacle}, \\
    \textit{ways}(i-1,j) + \textit{ways}(i,j-1) & \text{otherwise.}
  \end{cases}
\]
The base case also needs the obstacle check: $\textit{ways}(0,0) = 1$
\emph{unless} $(0,0)$ itself is an obstacle, in which case no path can
even begin, and the answer is $0$ overall.

\subsection*{Approach 1: Recursion}

\begin{lstlisting}
#include <bits/stdc++.h>
using namespace std;

int uniquePaths2Recursive(int i, int j, vector<vector<int>>& grid) {
    if (i < 0 || j < 0 || grid[i][j] == 1) return 0;
    if (i == 0 && j == 0) return 1;

    int fromAbove = uniquePaths2Recursive(i - 1, j, grid);
    int fromLeft  = uniquePaths2Recursive(i, j - 1, grid);
    return fromAbove + fromLeft;
}

int main() {
    vector<vector<int>> grid = {{0,0,0},{0,1,0},{0,0,0}};
    int m = grid.size(), n = grid[0].size();
    cout << uniquePaths2Recursive(m - 1, n - 1, grid) << endl; // 2
    return 0;
}
\end{lstlisting}

\begin{complexitybox}[Approach 1 Complexity]
\textbf{Time.} $O(2^{m+n})$.
\textbf{Space.} $O(m+n)$ recursion stack.
\end{complexitybox}

\subsection*{Approach 2: Memoization}

\begin{lstlisting}
#include <bits/stdc++.h>
using namespace std;

int uniquePaths2Memo(int i, int j, vector<vector<int>>& grid,
                      vector<vector<int>>& memo) {
    if (i < 0 || j < 0 || grid[i][j] == 1) return 0;
    if (i == 0 && j == 0) return 1;
    if (memo[i][j] != -1) return memo[i][j];

    int fromAbove = uniquePaths2Memo(i - 1, j, grid, memo);
    int fromLeft  = uniquePaths2Memo(i, j - 1, grid, memo);
    return memo[i][j] = fromAbove + fromLeft;
}

int main() {
    vector<vector<int>> grid = {{0,0,0},{0,1,0},{0,0,0}};
    int m = grid.size(), n = grid[0].size();
    vector<vector<int>> memo(m, vector<int>(n, -1));
    cout << uniquePaths2Memo(m - 1, n - 1, grid, memo) << endl; // 2
    return 0;
}
\end{lstlisting}

\begin{complexitybox}[Approach 2 Complexity]
\textbf{Time.} $O(m \cdot n)$.
\textbf{Space.} $O(m \cdot n)$ memo $+$ $O(m+n)$ recursion stack.
\end{complexitybox}

\subsection*{Approach 3: Tabulation}

\begin{lstlisting}
#include <bits/stdc++.h>
using namespace std;

int uniquePaths2Tabulation(vector<vector<int>>& grid) {
    int m = grid.size(), n = grid[0].size();
    vector<vector<int>> dp(m, vector<int>(n, 0));

    for (int i = 0; i < m; i++) {
        for (int j = 0; j < n; j++) {
            if (grid[i][j] == 1) {
                dp[i][j] = 0;
                continue;
            }
            if (i == 0 && j == 0) {
                dp[i][j] = 1;
                continue;
            }
            int fromAbove = (i > 0) ? dp[i-1][j] : 0;
            int fromLeft  = (j > 0) ? dp[i][j-1] : 0;
            dp[i][j] = fromAbove + fromLeft;
        }
    }
    return dp[m - 1][n - 1];
}

int main() {
    vector<vector<int>> grid = {{0,0,0},{0,1,0},{0,0,0}};
    cout << uniquePaths2Tabulation(grid) << endl; // 2
    return 0;
}
\end{lstlisting}

\subsection*{Dry run of the tabulation table}

\dryrun{Take $\textit{grid} = \begin{bmatrix}0&0&0\\0&1&0\\0&0&0\end{bmatrix}$.}

\begin{center}
\begin{tabular}{c|c|c|c}
 & $j=0$ & $j=1$ & $j=2$ \\ \hline
$i=0$ & 1 & 1 & 1 \\ \hline
$i=1$ & 1 & 0 (obstacle) & 1 \\ \hline
$i=2$ & 1 & 1 & 2
\end{tabular}
\end{center}

At $(1,1)$, the obstacle forces $\textit{dp}=0$ regardless of its
neighbors. At $(2,2)$: $\textit{dp}[1][2]+\textit{dp}[2][1] = 1+1=2$.
The answer is $2$: the two paths route either entirely along the top
then down the right edge, or entirely down the left edge then along the
bottom.

\begin{complexitybox}[Approach 3 Complexity]
\textbf{Time.} $O(m \cdot n)$.
\textbf{Space.} $O(m \cdot n)$ table, no recursion stack.
\end{complexitybox}

\subsection*{Approach 4: Space Optimization}

\begin{lstlisting}
#include <bits/stdc++.h>
using namespace std;

int uniquePaths2SpaceOptimized(vector<vector<int>>& grid) {
    int m = grid.size(), n = grid[0].size();
    vector<int> prevRow(n, 0);

    for (int i = 0; i < m; i++) {
        vector<int> currentRow(n, 0);
        for (int j = 0; j < n; j++) {
            if (grid[i][j] == 1) {
                currentRow[j] = 0;
                continue;
            }
            if (i == 0 && j == 0) {
                currentRow[j] = 1;
                continue;
            }
            int fromAbove = prevRow[j];
            int fromLeft  = (j > 0) ? currentRow[j-1] : 0;
            currentRow[j] = fromAbove + fromLeft;
        }
        prevRow = currentRow;
    }
    return prevRow[n - 1];
}

int main() {
    vector<vector<int>> grid = {{0,0,0},{0,1,0},{0,0,0}};
    cout << uniquePaths2SpaceOptimized(grid) << endl; // 2
    return 0;
}
\end{lstlisting}

\begin{complexitybox}[Approach 4 Complexity]
\textbf{Time.} $O(m \cdot n)$.
\textbf{Space.} $O(n)$ for the rolling row.
\end{complexitybox}

\begin{takeawaybox}
Adding a ``forbidden state'' constraint to an already-derived DP
recurrence is almost always a single extra case at the very top of the
recurrence (here, ``obstacle $\Rightarrow$ zero,'' checked before
anything else), with every other part of the four-stage derivation
unchanged. This is the same kind of small, composable extension already
seen when Combination Sum III added a length constraint to Combination
Sum II's base case.
\end{takeawaybox}
